\section{The \textsf{remove-copy} Algorithm}
\label{remove-copy}

We now study the \texttt{remove-copy} algorithm from ``ACSL By Example''.
This algorithm copies the contents from a source sequence to a destination
sequence whilst removing elements having a given value. The return value is the
length of the range.

\subsection{Adapting the Axiomatics}

The axiomatization provided in ``ACSL By Example'' is:

\listingname{remove-copy.axioms}
\cinput{removecopy/original.axioms}

The predicate \texttt{WhitherRemove} is defined by a \texttt{read} clause.
It must be adapted for the current version of the \textsf{WP} plug-in which
only implements a limited subset of \texttt{read} clauses:

\listingname{remove-copy.axioms [adapted]}
\cinput{removecopy/removecopy.axioms}

\subsection{Proving the Algorithm}

The specification of the \texttt{remove-copy} algorithm bases itself on the
\texttt{RemoveCopy} predicate, itself based on the \texttt{WhitherRemove}
function:

\listingname{removecopy.h}
\cinput{removecopy/removecopy.spec}

The implementation of the algorithm is:

\listingname{removecopy.c}
\cinput{removecopy/removecopy.impl}

\clearpage
The implementation can be partially proved correct against its specification by
running the \textsf{WP} plug-in: 
\fclogs{mutating}{removecopy}
